% ============================================================================
%  chapter5.tex
%  System Calls, File I/O, and Process Creation
% ============================================================================

\chapter{System Calls, File I/O, and Process Creation}

\section{The Operating System Is a Service Provider}

Everything before this chapter happened entirely inside a single
process: bytes in memory, structs on the heap, arrows between them.
That world is a comfortable one, but it is also a closed one. Real
software must reach \emph{outside} itself---to read a file, to spawn a
worker, to send a packet, to ask what time it is. Every one of those
acts crosses the boundary between your program and the operating
system.

An operating system is, at its core, a \emph{service provider with a
monopoly on the hardware}. Three jobs define it:

\begin{description}
  \item[Multiplexing] Many programs run at once on one CPU, one RAM,
    and a handful of disks. The OS interleaves their use of every
    device so that each program believes it has the machine to itself.
  \item[Protection] A program that misbehaves must not corrupt another
    program, the OS, or the hardware. The OS enforces this with
    hardware-enforced privilege levels.
  \item[Abstraction] The OS hides the terrifying variety of hardware---a
    dozen disk controllers, a hundred network cards---behind a small
    number of clean interfaces: files, processes, sockets, clocks.
\end{description}

\begin{verbatim}
   Without the OS, each program would have to:           With the OS:

   - drive the disk controller itself,                    read(fd, buf, n);
   - know the exact geometry of every disk,
   - police itself against corrupting other programs,
   - schedule its own CPU time,                          fork(); exec();
   - share the network card without colliding.           socket(); send();
\end{verbatim}

This chapter is the pivot of the book: we stop arranging memory inside
one process and start asking the OS for services. Every later chapter
is a deeper dive into one of these services. The chapter also teaches
the habit that makes systems programming survivable: \emph{every
OS call returns either success or an error, and a real program treats
both as normal control flow}.

\begin{osconcept}
An operating system exists precisely because hardware is shared,
dangerous, and diverse. Individual applications cannot be trusted to
multiplex the CPU or the disk among each other (they would starve and
corrupt one another), so resource management, protection, and uniform
device interfaces are problems that only a privileged, global piece of
software---the OS---can solve. That is why ``how do I read a file?''
is an \emph{OS} question and not a library question: the answer is a
contract with the kernel.
\end{osconcept}

\subsection{What the Kernel Is}

The word \emph{kernel} has been used casually here; let us pin it
down. The \emph{kernel} is the core of the operating system: the
privileged program that runs in kernel mode, owns the hardware, and
provides the system calls. When this book says ``the OS decides'' or
``the kernel enforces,'' it means this one program---the software
that controls the CPU, the memory, the disks, and the network card,
and that no other program may modify or even fully see.

\begin{verbatim}
   The kernel, drawn in the machine:

   +--------------------------------------------------------+
   |  THE KERNEL  -- one privileged program, in kernel mode  |
   |   owns: the CPU, memory pages, disks, network card      |
   |   provides: every system call you have used             |
   |   is: loaded at boot, resident for the machine's life   |
   +--------------------------------------------------------+
        ^                        ^
        | syscalls               | the hardware
        v                        v
   +--------------------------------------------------------+
   |  USER PROCESSES  (your programs, in user mode)          |
   |  many, untrusted, isolated from each other              |
   +--------------------------------------------------------+

   - "Kernel" is also the name of the actual binary: on Linux,
     /boot/vmlinuz (Chapter 15) IS the kernel, loaded at boot.
   - The kernel is the ONLY program allowed to touch the hardware.
     Everything your program does with the machine goes through it.
   - That is why every chapter's refrain--"only the kernel can do
     this"--is literally true: the CPU refuses to let anyone else.
\end{verbatim}

The kernel is not a distant concept; it is the direct addressee of
every system call in this book. \code{open}, \code{fork},
\code{mmap}, \code{poll}, and \code{connect} are all \emph{requests
to the kernel}. Everything since Chapter~5 has been a conversation
with one privileged program, and the rest of this book is learning
the terms of that conversation.

\section{User Mode and Kernel Mode}

To enforce protection, the CPU itself has at least two privilege
levels. In \emph{user mode}, your program runs. In \emph{kernel mode},
the OS runs. The difference is absolute: many instructions (turning off
interrupts, writing to device registers, remapping memory) simply
\emph{fault} if executed from user mode. The hardware forbids your
program from doing the dangerous things, so the OS does not have to
trust you.

\begin{verbatim}
   User mode (your program)              Kernel mode (the OS)

   +-----------------------------+       +--------------------------------+
   | main()                      |       | device drivers                |
   |   read(fd, buf, 4096);     |       | file systems (ext4, xfs)       |
   |   fork();                  |       | process scheduler             |
   |   write(1, msg, n);        |       | virtual memory manager        |
   |                            |       | open file table               |
   | your code can ONLY touch   |       | process table                 |
   | its OWN address space      |       |                              |
   +----------------------------+       | may touch ALL hardware        |
                |                       +--------------------------------+
                |  "please handle this for me"  (a system call)
                v
        +------------------+
        | privilege switch |
        | (CPU exception)  |
        +------------------+
\end{verbatim}

\begin{osconcept}
Protection cannot be a polite agreement between programs; it must be
\textbf{mechanically enforced}, and only the hardware can do that. By
having the CPU refuse dangerous instructions in user mode, the OS makes
misbehavior impossible-by-construction instead of just discouraged.
Deciding who gets kernel privilege, and keeping that boundary
leak-proof, is a core operating system problem.
\end{osconcept}

\section{System Calls: The Only Door}

The privileged operations are exposed to user programs through
\emph{system calls}: the tiny, fixed set of doors into the kernel.
Your C code rarely issues a system call directly; the C library
(\code{glibc}) provides wrappers that look like ordinary functions but
perform the privilege switch underneath.

\begin{verbatim}
   One system call, three layers:

   your code:            read(fd, buf, 4096);
        |
        v
   glibc wrapper:        place syscall number 0 and the arguments in
                         registers; execute the CPU instruction "syscall"
        |
        v
   kernel:               switch to kernel mode, look up syscall number 0,
                         validate the arguments, perform the operation,
                         store the result, switch back to user mode.

   A system call is NOT an ordinary function call. It changes privilege
   level, and the CPU does real work to make that safe. That is why
   system calls are comparatively expensive (hundreds of nanoseconds)
   and why programs do their best to do many file operations per call.
\end{verbatim}

\begin{verbatim}
   On x86-64 Linux, some system call numbers you are about to use:

   read      = 0        openat   = 257
   write     = 1        close    = 3
   exit      = 60       fork     = 57 (glibc wraps clone)
   lseek     = 8        execve   = 59
\end{verbatim}

Two properties of system calls matter enormously. First, they are
\emph{stable}: application code survives decades because the kernel
keeps syscall semantics constant. Second, they return
\emph{status}: almost every system call returns \code{-1} on failure
and records a reason in the variable \code{errno}. There is no
exception mechanism crossing the kernel boundary; there is a return
value, and you must read it.

The mechanical basis of the privilege switch is the \emph{interrupt}---
in this context also called a \emph{trap}. An interrupt is the CPU
abandoning its normal instruction stream to handle an exceptional
event: a key was pressed (a hardware interrupt), a page was missing
(Chapter~1), or a program executed the \code{syscall} instruction (a
software interrupt, the trap). Whatever the cause, the CPU saves its
state and jumps to a fixed handler \emph{in the kernel}. The word
``interrupt'' will return in Chapter~6 (signals interrupt your
program) and Chapter~10 (a system call interrupted by a signal): all
of them are the same CPU mechanism---``stop what you are doing, deal
with this, come back.''

\begin{osconcept}
The system call interface is the operating system's public API, and
designing it is an OS problem because it must be minimal (every entry
point is an attack surface and a maintenance contract), stable (users
depend on it for decades), and safe (the kernel must validate every
argument from an untrusted caller). Getting that boundary right is one
of the oldest and hardest problems in systems software.
\end{osconcept}

\section{Files: The Universal Abstraction}

The single most successful abstraction in computing is the \emph{file}:
a named sequence of bytes. Almost every device the OS owns is exposed
as a file---disks, terminals, pipes, even some clocks. The virtual file
system (VFS) lets the same \code{open}/\code{read}/\code{write} calls
work regardless of whether the bytes live on an ext4 disk, a network
share, or in another program's output.

\begin{verbatim}
   The file abstraction unifies everything:

     /dev/tty     -> a terminal            (read/write characters)
     /proc/cpuinfo-> data manufactured on the fly by the kernel
     fifo          -> a pipe between two programs
     socket        -> network data
     /home/you/f  -> bytes on a disk

   From your program's point of view each is "a sequence of bytes you
   can read or write." The OS sorts out where the bytes really live.
\end{verbatim}

\begin{osconcept}
Naming, layout, and storage of persistent data are an OS problem
because only the kernel can give programs a \emph{namespace} (the file
tree) that is shared, protected by permissions, and backed by arbitrary
physical devices. A library cannot create a file; it can only call the
syscall that does. Files are the OS's way of saying ``here is a uniform
interface to everything persistent or remote.''
\end{osconcept}

\section{File Descriptors: Your Handles on the Kernel}

When you \code{open} a file, the kernel creates an entry in its own
bookkeeping---tracking the file position, the access mode, the
permissions---and hands you back a small integer, the \emph{file
descriptor}. From then on you never name the path again; you name the
descriptor. The kernel keeps the path lookup and the caching behind
your back.

\begin{lstlisting}[caption={Opening a file for reading.}]
#include <stdio.h>
#include <stdlib.h>
#include <fcntl.h>
#include <unistd.h>

int main(void) {
    int fd;

    fd = open("hello.txt", O_RDONLY);
    if (fd < 0) {
        perror("open");
        exit(1);
    }

    printf("opened file, fd = %d\n", fd);

    close(fd);     /* release the kernel's entry */
    return 0;
}
\end{lstlisting}

\begin{verbatim}
   Your process                          The kernel
   +-------------------------+            +---------------------------+
   | file descriptor table   |            | open file table           |
   |  0  ->  stdin      ----+----------->|  position: 0              |
   |  1  ->  stdout     ----+----------->|  mode: O_RDONLY           |
   |  2  ->  stderr     ----+----------->|  -> inode (disk file)     |
   |  3  ->  hello.txt  ----+----------->|                           |
   |  4  ->  (free)         |            |                           |
   +-------------------------+            +---------------------------+

   - Descriptors 0, 1, 2 are ALWAYS present: stdin, stdout, stderr.
   - open() returns the LOWEST free descriptor (here 3).
   - close(3) removes the entry; the number is then reusable.
   - The path "hello.txt" is resolved ONCE at open(); everything after
     uses the number.
\end{verbatim}

The flags matter. \code{O\_RDONLY} opens read-only; \code{O\_WRONLY}
write-only; \code{O\_RDWR} both. \code{O\_CREAT} says ``create the file
if it does not exist'' and then requires a third argument, the mode
bits. \code{O\_TRUNC} empties an existing file first. \code{O\_APPEND}
forces every write to go to the end, which is how logs stay logs even
when several processes write to them.

\begin{osconcept}
A file descriptor is a \emph{capability}: a small token that says
``this process may access that kernel object.'' Managing the per-process
descriptor table, the shared open-file table, and the reference
counting that decides when a file is truly closed is pure OS work.
Your program sees an integer; the kernel sees a reference into
structures it alone can touch. That indirection is what allows
permissions to be checked once at open time and then enforced cheaply
for every subsequent call.
\end{osconcept}

\begin{researchnote}[The paper that invented all of this]
Every abstraction in this chapter---files as byte streams, descriptors,
pipes, processes---was designed in one place, in 1974, and written up
in the most influential systems paper ever published: ``The UNIX
Time-Sharing System,'' by Ken Thompson and Dennis Ritchie. If you read
one paper from this entire book, read this one.

\begin{verbatim}
   An extract (a real sentence from the paper):

       "The most important job of UNIX is to provide a file system.
        From the point of view of the user, there are three kinds of
        files: ordinary disk files, directories, and special files."

   What this emphasis on the file system means, in our chapter's terms:
   the authors chose to make nearly EVERYTHING a file -- ordinary data,
   directories, even devices -- so that the tools for reading and
   writing one thing work for all of them. That single decision is why
   our open/read/write pattern is so uniform, and why descriptors 0,
   1, and 2 exist at all.
\end{verbatim}

Thompson and Ritchie published it after years of internal use at Bell
Labs, and the paper is famous for being both short and astonishingly
clear. Reading it is a masterclass in explaining an operating system
in prose---exactly the style this book tries to imitate.

\paperref{K.~Thompson, D.~M.~Ritchie, ``The UNIX Time-Sharing System,''
Communications of the ACM 17(7), 1974}{https://www.cs.princeton.edu/courses/archive/spring14/cos217/reading/unix.pdf}
\end{researchnote}

\section{Reading and Writing}

\subsection{What Is a Buffer?}

You will meet the word \emph{buffer} in every chapter of this book, so
let us fix its meaning once and for all. A \emph{buffer} is a
contiguous region of memory used to hold data temporarily while it is
in transit between two places. The key word is \emph{temporarily}: a
buffer is a staging area, not a destination. Data flows in, sits
briefly, and flows out.

Three facts of real hardware and real kernels make buffers
unavoidable in systems programming:

\begin{itemize}
  \item Devices deliver data in \emph{chunks}, never one byte at a
    time on demand. A disk reads a whole sector; a network card
    delivers a whole packet. A buffer is where a chunk waits while
    your code catches up.
  \item Crossing the kernel boundary is expensive. Doing I/O in big
    batches through a buffer costs far less than many tiny transfers.
  \item Producers and consumers run at different speeds. A buffer is
    the shock absorber between them.
\end{itemize}

\begin{verbatim}
   The central picture: data in transit through a staging area.

      producer -----> [  B U F F E R  ] -----> consumer
                       (contiguous bytes
                        waiting their turn)

   In the copy_file() program below:
      read()  fills buf from the disk   (producer: the kernel)
      write() drains buf to the disk    (consumer: the kernel)
      your buffer sits between the two.
\end{verbatim}

In C a buffer is simply an array---or a block from \code{malloc}.
Our copy loop's \code{char buf[CHUNK]} is a buffer of 4096 bytes:
\code{read} deposits bytes into it, \code{write} drains them, and the
two operations alternate until the file is exhausted. Notice that the
buffer never grows; it is reused, which is the whole point of a
staging area.

\begin{verbatim}
   Where buffers appear throughout this book:

     file I/O      :  char buf[4096];  this chapter
     strings       :  strncpy targets, flexible array tails (ch.3)
     the kernel    :  the buffer cache -- disk blocks cached in RAM
     pipes         :  the kernel's pipe buffer  (ch.7)
     networking    :  a socket's receive/send buffers (ch.17)
     producer/consumer : a RING buffer (circular queue) (ch.11)

   They are all the same idea: contiguous bytes held in transit.
\end{verbatim}

The magnitude of the win is worth stating with numbers. A system call
costs hundreds of nanoseconds because of the privilege switch we saw
earlier in this chapter.

\begin{verbatim}
   Reading a 1 GB file:

     one byte at a time:   read(fd, &b, 1);      ~1,000,000,000 syscalls
     4096-byte chunks:     read(fd, buf, 4096);    ~250,000      syscalls

     The buffer turns a billion expensive crossings of the kernel
     boundary into a quarter-million. That is the difference between a
     program that finishes in seconds and one that takes an hour.
\end{verbatim}

\begin{osconcept}
Buffering is not just a user-mode habit; the kernel buffers too. The
\emph{buffer cache} keeps recently-read disk blocks in RAM so that a
second read of the same block never touches the disk again, and write
coalescing lets the kernel schedule big sequential writes instead of
thousands of tiny ones. Deciding how much to buffer, when to flush,
and how to share the cache fairly among processes is an operating
system problem. Your user-space buffer is the interface to that
machinery.
\end{osconcept}

\subsection{The Partial I/O Contract}

\code{read(fd, buf, n)} asks the kernel to copy up to \code{n} bytes
from the file into your buffer; it returns how many it actually copied.
\code{write(fd, buf, n)} does the reverse. The subtlety that beginners
miss, and that makes systems programming distinct from scripting, is
that these calls may transfer \emph{fewer bytes than requested}, and
that is not an error---it is the normal contract of an I/O device.

\begin{lstlisting}[caption={A correct copy loop: reads and writes in chunks.}]
#include <stdio.h>
#include <stdlib.h>
#include <fcntl.h>
#include <unistd.h>

#define CHUNK 4096

/* Copy src into dst. Returns 0 on success, -1 on failure. */
int copy_file(const char *src, const char *dst) {
    int in_fd;
    int out_fd;
    char buf[CHUNK];
    ssize_t got;
    ssize_t written;
    ssize_t r;
    int failed;

    in_fd = open(src, O_RDONLY);
    if (in_fd < 0) {
        perror("open(src)");
        return -1;
    }

    out_fd = open(dst, O_WRONLY | O_CREAT | O_TRUNC, 0644);
    if (out_fd < 0) {
        perror("open(dst)");
        close(in_fd);
        return -1;
    }

    failed = 0;
    while (1) {
        got = read(in_fd, buf, sizeof(buf));
        if (got < 0) {
            perror("read");
            failed = 1;
            break;
        }
        if (got == 0) {
            break;          /* end of file: nothing left to copy */
        }

        /* write() may transfer fewer bytes than we asked; loop until
           every byte we read has been written.                    */
        written = 0;
        while (written < got) {
            r = write(out_fd, buf + written, (size_t) (got - written));
            if (r < 0) {
                perror("write");
                failed = 1;
                break;
            }
            written = written + r;
        }
        if (failed) {
            break;
        }
    }

    close(in_fd);
    close(out_fd);
    return failed ? -1 : 0;
}

int main(int argc, char *argv[]) {
    if (argc != 3) {
        fprintf(stderr, "usage: %s src dst\n", argv[0]);
        return 1;
    }
    return copy_file(argv[1], argv[2]) == 0 ? 0 : 1;
}
\end{lstlisting}

\begin{verbatim}
   Why the inner loop?  Because the contract says "at most n bytes":

   write(out_fd, buf + written, got - written);
                                   ^^^^^^^^^^^
   may write 1, or 10, or all 4096. The loop keeps calling until the
   count is exhausted. For files on disk this rarely matters; for
   pipes, terminals, and network sockets it matters constantly. Build
   the loop once, copy it into every program you write.
\end{verbatim}

\begin{verbatim}
   The read contract:

     read(fd, buf, n)  =  n          n bytes copied, exactly as asked
     read(fd, buf, n)  =  k < n      k bytes copied (partial); try again
     read(fd, buf, n)  =  0          end of file, no more bytes
     read(fd, buf, n)  =  -1         error; inspect errno
\end{verbatim}

Two design choices in this program are deliberate. The buffer is
\code{4096} bytes---one page---because that matches the granularity the
kernel and the disk like. And we check every return value, because in
systems code the ``impossible'' failure is not impossible; it is just
waiting for the right day to happen.

\begin{osconcept}
I/O devices are asynchronous, shared, and slow, and the kernel sits
between your program and all of that: it decides when a read can
complete (the device may still be spinning), it buffers and caches so
repeated small reads do not hammer the disk, and it interleaves
requests from all processes fairly. The partial-I/O contract, the
buffer cache, and the scheduling of device requests are therefore
operating system problems; your program merely states what it wants and
pays attention to what it got.
\end{osconcept}

\section{Seeking: Moving the File Position}

Each open file carries a \emph{position}: the offset where the next
read or write will happen. Reads and writes advance it. \code{lseek}
moves it without touching any data---seeking past the end then writing
creates a \emph{hole} of zero bytes, which is how sparse files are born.

\begin{lstlisting}[caption={Finding a file's size by seeking to its end.}]
#include <stdio.h>
#include <fcntl.h>
#include <unistd.h>
#include <sys/stat.h>

/* Return the size in bytes of the file opened as fd, or -1. */
off_t file_size(int fd) {
    off_t end;

    end = lseek(fd, 0, SEEK_END);
    if (end < 0) {
        perror("lseek");
        return -1;
    }
    return end;
}
\end{lstlisting}

\begin{verbatim}
   A file with a hole:

        logical bytes:   0       1024     2048    4096
                         +--------+--------+--------+
                         | 'A'*1k | (hole) | 'B'*2k |
                         +--------+--------+--------+
   physical disk blocks: allocated  nothing  allocated
                         (the hole costs no disk space; reads of the
                          hole return zero bytes)

   lseek(fd, 2048, SEEK_SET);  then write(fd, buf, 2048);
   fills the hole only if you actually write into it.
\end{verbatim}

\begin{osconcept}
The file position is kernel state attached to an open file, not to the
file on disk or to your program. When two processes open the same file,
each gets its own position, while \code{O\_APPEND} deliberately makes
the kernel bump the position before every write so concurrent appenders
never clobber each other. Deciding where the position lives, and how
concurrent access to it is serialized, is operating system work.
\end{osconcept}

\section{errno: The Error Channel}

Every system call that can fail records the reason in a thread-local
integer named \code{errno}. Two functions turn it into human words:
\code{perror(prefix)} prints \code{prefix} plus the message, and
\code{strerror(errno)} returns the message string alone. The discipline
is to check the failure and report \emph{immediately}, while the
meaning is still fresh.

\begin{lstlisting}[caption={The error-handling pattern used everywhere in this book.}]
#include <stdio.h>
#include <stdlib.h>
#include <fcntl.h>
#include <errno.h>

int open_or_die(const char *path) {
    int fd;

    fd = open(path, O_RDONLY);
    if (fd < 0) {
        fprintf(stderr, "cannot open %s: %s\n",
                path, strerror(errno));
        exit(1);
    }
    return fd;
}
\end{lstlisting}

\begin{verbatim}
   What each failure actually means:

     open("f.txt", O_RDONLY)  ->  -1, errno = ENOENT   No such file
     open("/root/x", O_RDONLY)->  -1, errno = EACCES   Permission denied
     open("d/", O_RDONLY)     ->  -1, errno = EISDIR   Is a directory
     read(fd, ...)            ->  -1, errno = EINTR    Interrupted by signal

   perror("open")  prints:    open: No such file or directory
\end{verbatim}

\begin{ruleofthumb}[Report errors where they happen]
Do not swallow errors, and do not report them a mile from the cause.
The instant a call returns \code{-1}, print or log it with enough
context (\code{errno} and the operation) and decide the failure policy
right there: retry, degrade, or abort. Programs that defer error
handling are the ones that ship the confusing bug reports.
\end{ruleofthumb}

\section{Standard Streams and Redirection}

Descriptors 0, 1, and 2 are already open: standard input, standard
output, and standard error. \code{printf} writes to descriptor 1;
\code{fprintf(stderr, ...)} to descriptor 2. The shell's
\code{>}, \code{<}, and \code{2>} operators are nothing but the shell
rearranging what those descriptors point at before it \code{exec}s your
program.

\begin{verbatim}
   What the shell does for "myprog > out.txt 2> err.txt":

     1. open("out.txt", O_WRONLY|O_CREAT|O_TRUNC, 0644) -> fd 3
     2. open("err.txt", O_WRONLY|O_CREAT|O_TRUNC, 0644) -> fd 4
     3. dup2(3, 1);   /* make descriptor 1 ALSO name out.txt  */
     4. dup2(4, 2);   /* make descriptor 2 ALSO name err.txt  */
     5. close(3); close(4);
     6. exec("myprog");

   myprog sees stdout as out.txt and stderr as err.txt, with no
   knowledge that a redirection ever happened. That is the power of
   the descriptor indirection: the program does not care what its
   streams are connected to.
\end{verbatim}

\begin{lstlisting}[caption={Redirecting stdout inside your own program.}]
#include <stdio.h>
#include <stdlib.h>
#include <fcntl.h>
#include <unistd.h>

int main(void) {
    int fd;

    fd = open("out.txt", O_WRONLY | O_CREAT | O_TRUNC, 0644);
    if (fd < 0) {
        perror("open");
        exit(1);
    }

    /* Make descriptor 1 (stdout) name the same file as fd. */
    if (dup2(fd, 1) < 0) {
        perror("dup2");
        exit(1);
    }
    close(fd);       /* we no longer need the original descriptor */

    printf("this line lands in out.txt\n");
    return 0;
}
\end{lstlisting}

\code{dup2(fd, 1)} is the magic word: it copies the kernel object that
\code{fd} names onto descriptor 1, atomically. This single call is the
foundation of every shell redirection, every pipeline, and most of the
inter-process communication we will build in Chapter~7.

\begin{osconcept}
Redirection works because the kernel, not the program, owns what each
descriptor names. A program cannot tell whether descriptor 1 is a
terminal, a file, or a pipe---and it does not need to. Making the
connection between a descriptor number and the underlying object a
changeable piece of kernel state, mutable by \code{dup2}, is what lets
an external agent (the shell) rewire a program's world from outside.
That is a subtle, deliberate piece of OS design.
\end{osconcept}

\section{Processes: What the OS Schedules}

A \emph{process} is a running program: one virtual address space, one
set of registers, one execution thread, and a set of kernel resources
(open files, a working directory, environment variables). The OS keeps
one record per process---the process control block (PCB)---and the
\emph{scheduler} decides, thousands of times per second, which process
gets the CPU next. From the OS's point of view, a process is a unit of
accounting and isolation; from your program's point of view, it is
``everything I am right now.''

\begin{verbatim}
   The OS keeps MANY processes in flight on a few CPUs:

     CPU 0:  [ A ] [ C ] [ A ] [ D ] ...    (each column a time slice)
     CPU 1:  [ B ] [ A ] [ D ] [ B ] ...

   Each process THINKS it owns the machine; the scheduler interleaves
   them. Switching is fast but not free: registers, the memory map,
   and the page tables must be swapped ("context switch").

   Isolation: A cannot read or write B's memory. The OS gives each
   process its own private address space:

     Process A               Process B
     +---------+             +---------+
     | stack   |             | stack   |
     | heap    |             | heap    |
     | data    |             | data    |
     | text    |             | text    |
     +---------+             +---------+
       \                       /
        \  shared ONLY through explicit OS channels: files, pipes,
         sockets -- never by poking another process's memory.
\end{verbatim}

\begin{osconcept}
Multiplexing the CPU is the quintessential OS problem. The hardware has
one instruction stream per core; the OS manufactures the illusion of
concurrency by switching processes, and it must do so \emph{fairly}
(no process starves), \emph{safely} (one process cannot corrupt
another's memory), and \emph{cheaply} (switches are not too frequent).
Without the OS, cooperating programs could not share a CPU at all.
\end{osconcept}

\section{fork(): The Process Duplication Primitive}

How are new processes created? On Unix, the primitive is \code{fork}:
the kernel makes a near-exact copy of the calling process. The child is
a new process with its own copy of the parent's address space, its own
copy of the descriptor table, and its own process id. After the call,
\emph{two} processes are running the same code from the same line.

The cleverness of the implementation is that the copy is nearly free at
first. The parent's pages are shared with the child and marked
\emph{copy-on-write} (COW): the physical pages are only duplicated when
one of the two processes actually writes to them. Most programs
\code{fork} and then immediately \code{exec}, so most pages are never
copied at all.

\begin{lstlisting}[caption={fork returns twice: once in each process.}]
#include <stdio.h>
#include <stdlib.h>
#include <unistd.h>
#include <sys/types.h>

int main(void) {
    pid_t pid;

    pid = fork();
    if (pid < 0) {
        perror("fork");
        exit(1);
    }

    if (pid == 0) {
        /* This branch runs ONLY in the child. */
        printf("child:  my pid = %d, my parent's pid = %d\n",
               getpid(), getppid());
        exit(0);
    }

    /* This branch runs ONLY in the parent. */
    printf("parent: my pid = %d, my child's pid = %d\n",
           getpid(), pid);

    return 0;
}
\end{lstlisting}

\begin{verbatim}
   One call, two worlds:

                              fork()
                 parent -----------------> parent (returns child's pid)
                    |   |                     (the caller CONTINUES here)
                    |   +---- child ---------> child (returns 0)
                    |       (the child starts here, at the SAME line,
                    |        with a COPY of the parent's memory)

   The two processes are now independent: the child can change its
   copy of a variable and the parent's copy is untouched, thanks to
   copy-on-write the change happens on a page that is now private.
\end{verbatim}

Two facts about the return value drive every use of fork:

\begin{verbatim}
   In the PARENT, fork returns the child's pid (a positive number).
   In the CHILD,  fork returns 0.
   On failure,    fork returns -1 and no child is created.

   The single branch -- if (pid == 0) child else parent -- is the
   universal shape. After fork, parent and child have DIFFERENT pids,
   DIFFERENT copies of memory, but the SAME open files (the descriptor
   table is copied, so both can write to fd 1, for example).
\end{verbatim}

\begin{osconcept}
Creating a process is a privileged act: the kernel must allocate a new
process control block, a new address space, new page tables, and a new
scheduler entry, then wire the child into the process tree. A user
program cannot do any of that. The elegance of \code{fork} is that it
makes process creation almost free by deferring all copying via
copy-on-write---a decision only the OS can make, because only the OS
controls the page tables.
\end{osconcept}

\subsection{vfork: The Half-Fork}

There is an older, rawer sibling of \code{fork} that you will meet in
old code and in some fast libraries, and it deserves a name:
\code{vfork} is the \emph{half-fork}. Where \code{fork} gives the
child a copy (or a copy-on-write view) of the address space,
\code{vfork} gives the child \emph{the parent's very own pages}, and
\emph{suspends the parent} until the child either \code{exec}s or
\code{\_exit}s.

\begin{verbatim}
   fork  :  parent and child BOTH run, each with (eventually) its own
            copy of memory. The child can do anything the parent can.

   vfork :  only the child runs; the parent is FROZEN. The child is
            using the parent's memory directly -- no copy at all --
            and may ONLY call exec() or _exit(). That is the entire
            deal.

   Why it existed: before copy-on-write was invented, fork() copied
   the whole address space. For the overwhelmingly common
   "fork, then immediately exec a new program," that copy was pure
   waste. vfork skipped it: no copy, no COW, no memory at all -- the
   child borrows the parent's address space for the instant it needs
   it, then execs and leaves it behind.

   Why it is dangerous: because the child runs on the PARENT'S stack
   and memory, it must not modify anything, must not call ordinary
   functions (malloc, printf), and must never return from its own
   main -- if it does, it corrupts the suspended parent. The POSIX
   rules are absolute: in the vfork child you may call only
   async-signal-safe functions, exec, or _exit.
\end{verbatim}

On modern systems, \code{fork} is copy-on-write and costs almost
nothing, so \code{vfork} has lost most of its reason to exist. POSIX
keeps it for rare, performance-critical fork-and-exec paths, and the
standard is explicit about what the vfork child may do: call only
async-signal-safe functions, \code{exec}, or \code{\_exit}. The
half-fork survives in tight code that wants the absolute minimum
between two processes---and it is worth recognizing on sight, because
its constraints are the opposite of \code{fork}'s freedom.

\begin{ruleofthumb}[Never write new code with vfork]
Use \code{fork}. The copy-on-write \code{fork} of any modern kernel
is cheap, and it does not freeze the parent or forbid ordinary
library calls in the child. If you ever meet \code{vfork} in the
wild, read the word as a warning sign: the code around it is older
and stricter, and the child may only \code{exec} or \code{\_exit}.
\end{ruleofthumb}

\section{exec(): Replacing the Process Image}

\code{fork} duplicates the process; \code{exec} \emph{replaces} it.
The \code{exec} family (\code{execl}, \code{execv}, \code{execvp},
\code{execve}) loads a new program file into the current process,
discarding the old address space and starting the new program at its
\code{main}. The process id does not change. If \code{exec} succeeds,
it never returns---there is no old program left to return to. If it
fails, it returns \code{-1} and you are still in the old program.

\begin{lstlisting}[caption={Fork + exec: the standard recipe for "run this program".}]
#include <stdio.h>
#include <stdlib.h>
#include <unistd.h>
#include <sys/types.h>
#include <sys/wait.h>

/* Run the program named by argv[0] with arguments argv[1..].
   Returns the child's exit status, or -1 on failure. */
int run_command(char *const argv[]) {
    pid_t pid;
    int status;

    pid = fork();
    if (pid < 0) {
        perror("fork");
        return -1;
    }

    if (pid == 0) {
        /* CHILD: become the new program. execvp searches the PATH
           for argv[0]. It only returns here if it FAILED.       */
        execvp(argv[0], argv);
        perror("execvp");
        exit(127);
    }

    /* PARENT: wait for the child to finish, then report its status. */
    if (waitpid(pid, &status, 0) < 0) {
        perror("waitpid");
        return -1;
    }
    if (WIFEXITED(status)) {
        return WEXITSTATUS(status);
    }
    return -1;
}

int main(int argc, char *argv[]) {
    if (argc < 2) {
        fprintf(stderr, "usage: %s command [args...]\n", argv[0]);
        return 1;
    }
    return run_command(argv + 1);
}
\end{lstlisting}

\begin{verbatim}
   Life of a process that runs "ls -l":

   main()                       (your program)
     fork() --> [child]
     parent: waitpid()          (sits, waiting)
     child:  execvp("ls", ...)
             (the child's address space is REPLACED by /bin/ls;
              the pid is UNCHANGED)
     /bin/ls runs, prints, exits with status 0
   parent: waitpid() returns 0  (the child's exit status)
\end{verbatim}

\begin{verbatim}
   fork + exec = "build a new shell of a process, then fill it":
     fork creates a private address space with your image;
     exec destroys that image and loads a new one.
   The two steps are kept separate so the child can do setup between
   them -- close a descriptor, dup2 a pipe, change directories --
   which is exactly how shells and daemons are built.
\end{verbatim}

\begin{osconcept}
Loading an executable image is privileged work: the kernel must parse
the file format, map the code and data segments into a fresh address
space, set up the stack and the \code{argv}/\code{envp} arrays, and
honor the file's permissions and security flags. No user-mode code can
legitimately conjure a running program from a file---only the kernel
can. That \code{exec} preserves the pid and the open descriptors while
replacing the memory is a carefully chosen OS semantic that makes
processes composable.
\end{osconcept}

\section{wait() and waitpid(): Reaping Children}

When a child exits, it does not vanish. It enters the \emph{zombie}
state: all of its resources are released except a tiny record holding
its pid and its exit status. That record exists so the parent can
discover \emph{how} the child died. The parent must call
\code{wait} or \code{waitpid} to collect the record, at which point
the zombie is reaped and the pid becomes reusable.

\begin{verbatim}
   The life cycle of a child:

      fork() ---> running ---> exiting ---> ZOMBIE ---> reaped
                        |                      |
                        +------ exits --------++------ wait() collects it

   - A ZOMBIE consumes almost no resources, but its pid is taken.
   - If the parent NEVER waits, the zombie lingers until the parent
     exits; then init (pid 1) adopts it and reaps it.
   - If the parent exits FIRST, the child becomes an ORPHAN and is
     adopted by init (pid 1), which reaps it when it finally dies.

   The rule: every fork() deserves a matching wait(). Forgetting the
   wait leaks a zombie per child -- a classic "my server grew to
   thousands of processes" bug.
\end{verbatim}

\begin{lstlisting}[caption={waitpid with the exit-status macros.}]
#include <stdio.h>
#include <stdlib.h>
#include <unistd.h>
#include <sys/types.h>
#include <sys/wait.h>

int main(void) {
    pid_t pid;
    int status;

    pid = fork();
    if (pid < 0) {
        perror("fork");
        exit(1);
    }

    if (pid == 0) {
        printf("child exiting with status 3\n");
        exit(3);
    }

    if (waitpid(pid, &status, 0) < 0) {
        perror("waitpid");
        exit(1);
    }

    if (WIFEXITED(status)) {
        printf("child exited with status %d\n", WEXITSTATUS(status));
    } else if (WIFSIGNALED(status)) {
        printf("child killed by signal %d\n", WTERMSIG(status));
    }

    return 0;
}
\end{lstlisting}

\begin{verbatim}
   The status integer packs three stories:
     WIFEXITED(s)    : did it exit normally?
     WEXITSTATUS(s)  : if so, its exit code (0..255)
     WIFSIGNALED(s)  : was it killed by a signal?
     WTERMSIG(s)     : which signal?

   Without wait(), the parent has NO way to know these. And because
   a signal can kill a process at ANY moment, a parent that ignores
   wait() cannot even be sure its children succeeded.
\end{verbatim}

\begin{osconcept}
Reaping children is an invariant the OS maintains for the whole
process tree: no pid is ever reused while a parent might still ask
about it, and no orphan is ever left uncollected. Keeping the exit
status available until the parent explicitly retrieves it, adopting
orphans, and recycling pids safely are all kernel bookkeeping
decisions---your program just calls \code{waitpid}, and the OS handles
the rest. The zombie rule exists so the answer to ``what happened to my
child?'' is never lost.
\end{osconcept}

\section{Practical: A Tiny Command Runner}

Let us combine everything from this chapter into one tool: a miniature
shell for a single command. It forks, lets the child \code{exec} the
requested program, waits for it, and prints how it ended. This program
is a few dozen lines, yet it is the beating heart of every shell,
service manager, and job-control system ever written.

\begin{lstlisting}[caption={runone.c: fork, exec, wait -- the whole chapter at once.}]
#include <stdio.h>
#include <stdlib.h>
#include <unistd.h>
#include <sys/types.h>
#include <sys/wait.h>

int main(int argc, char *argv[]) {
    pid_t pid;
    int status;

    if (argc < 2) {
        fprintf(stderr, "usage: %s command [args...]\n", argv[0]);
        return 1;
    }

    pid = fork();
    if (pid < 0) {
        perror("fork");
        return 1;
    }

    if (pid == 0) {
        /* Child: become the requested program. */
        execvp(argv[1], argv + 1);
        perror("execvp");
        exit(127);          /* conventional "command not found" status */
    }

    /* Parent: wait for the child and report the outcome. */
    if (waitpid(pid, &status, 0) < 0) {
        perror("waitpid");
        return 1;
    }

    if (WIFEXITED(status)) {
        printf("command exited with status %d\n", WEXITSTATUS(status));
    } else if (WIFSIGNALED(status)) {
        printf("command killed by signal %d\n", WTERMSIG(status));
    }

    return 0;
}
\end{lstlisting}

\begin{verbatim}
   Running it:

     $ ./runone echo hello world
     hello world
     command exited with status 0

     $ ./runone ls /nonexistent
     ls: cannot access '/nonexistent': No such file or directory
     command exited with status 2

     $ ./runone /bin/sh -c 'kill -9 $$'
     command killed by signal 9

   One small program demonstrates: file descriptors survive exec
   (echo used the same fd 1 we inherited), argv is rebuilt by the
   kernel, and wait() reveals whether the child exited or was killed.
\end{verbatim}

The next time you type a command into a shell, know that you are
watching this chapter run: the shell \code{fork}s, you are the child
mid-\code{exec}, and the prompt you get back is the shell's
\code{waitpid} returning. Everything we built with structs and pointers
is what runs \emph{inside} such a process; this chapter is what creates
it.

\section{Structures in Bytes}

The byte structures this chapter touched are the tables the kernel
builds and the stack it prepares at \code{exec} time.

\begin{verbatim}
   The argv array on the new process's stack (high to low):

   +-------------------------------+
   | "two words\0"     argv[2]     |
   | "alpha\0"         argv[1]     |   the strings, each NUL-terminated
   | "./prog\0"        argv[0]     |
   +-------------------------------+
   | NULL                          |   argv[argc] == NULL
   | argv[2]  (pointer, 8B)        |
   | argv[1]  (pointer, 8B)        |   an array of 8-byte pointers
   | argv[0]  (pointer, 8B)        |
   +-------------------------------+
   | argc (8B)                     |   the count
   +-------------------------------+

   The kernel's descriptor table is an array of pointers to file
   objects; each file object carries the position, mode, and flags,
   and points at the inode. dup2() copies a pointer, and close()
   drops a reference count.
\end{verbatim}

A stack frame has the same shape as Chapter~1's: an 8-byte return
address pushed by \code{call}, the saved frame pointer, then the
locals. System calls add no byte structures to your process; they are
requests that the kernel's own records---descriptor table, open-file
table, process table---do the work.

\section{Summary}

The chapter's arc, in one picture:

\begin{verbatim}
   your program (user mode)
   +------------------------------------------------------------+
   | structs & pointers   (chapters 1-3)                        |
   | library calls: malloc, printf, strcmp                       |
   +----------------------------+-------------------------------+
                                | system calls (the only door)
                                v
   the kernel (kernel mode)  +---------------------------------+
   | open/read/write/lseek/close  -> files                     |
   | fork/exec/waitpid           -> processes                  |
   | dup2                        -> redirection & wiring      |
   | descriptor table, open file table, scheduler, zombies    |
   +-----------------------------------------------------------+

   discipline: every call is checked, every error is reported with
   errno, every fork is matched with a wait, every fd is closed.
\end{verbatim}

\begin{osconcept}
This entire chapter is ``the OS problem'' in miniature: a safe,
stable, shared interface between untrusted user code and privileged
hardware. The abstractions you met---files, descriptors, processes,
wait--exist so that many programs can coexist safely on one machine.
Understanding that these are \emph{services with contracts}, not
helper functions, is what separates a C programmer from a systems
programmer.
\end{osconcept}

In the next chapter we meet the events that can interrupt a process at
any moment---signals---and learn how to handle, ignore, and block
them without corrupting our own state.

